\documentclass[10pt,a4paper]{amsart}
\usepackage[margin=19mm]{geometry}
\usepackage{amsmath,amssymb,mathtools}
\usepackage[scr=rsfs,cal=euler]{mathalfa}
\usepackage{xfrac}
\usepackage[unicode,hidelinks,bookmarks=false]{hyperref}
\newcommand{\ie}{i.e.\ }
\newcommand{\cf}{cf.\ }
\newcommand{\resp}{respectively\ }
\newcommand{\Subsection}{\subsection}
\providecommand{\nobreakdash}{\nobreak\hbox{-}}
\setlength{\emergencystretch}{2em}
\multlinegap0pt
\usepackage{bm}
\usepackage{bbm}
\usepackage{dsfont}
\usepackage{xspace}
\usepackage{cancel}
\usepackage{mathtools}
\usepackage{amsthm}
\makeatletter
\@ifundefined{theorem}{\newtheorem{theorem}{Theorem}}{}
\@ifundefined{corollary}{\newtheorem{corollary}[theorem]{Corollary}}{}
\@ifundefined{lemma}{\newtheorem{lemma}[theorem]{Lemma}}{}
\@ifundefined{proposition}{\newtheorem{proposition}[theorem]{Proposition}}{}
\makeatother
\title[Relative ultragraph algebras and infinite interval maps]{Relative ultragraph algebras and infinite interval maps}
\author{Ben-Hur Eidt, Daniel Gon\c{c}alves, Danilo Royer}
\date{Source: arXiv:2508.19835v1 (CC BY 4.0)}
\begin{document}
\maketitle

\noindent\emph{Source and licence.} Relative ultragraph algebras and infinite interval maps. Version arXiv:2508.19835v1 (CC BY 4.0), distributed under the Creative Commons Attribution 4.0 International licence, as recorded in the arXiv licence metadata for this exact version and shown on the abstract page https://arxiv.org/abs/2508.19835v1. Full rights notice: licenses/arxiv-2508.19835-CC-BY-4.0.txt.\par\smallskip

\noindent\textit{Editorial scope.} Counted unit: the Proposition whose source label is \texttt{isomorphismrelative} in the article (source0.tex lines 823--829, source label \texttt{isomorphismrelative}), delivered together with the article's own complete original proof of it (source0.tex lines 830--1166, one \texttt{proof} environment of 337 lines). No mathematics is written, repaired or re-derived: statement and proof are the article's own text line for line. Required context retained uncounted: esfriou, vai dar tainha (lines 424--443); decomposition2 (lines 688--696); contas1 (lines 760--776). Every definition, display and label of those lines is retained unchanged. Omitted from this package and not redistributed: 1--423, 444--687, 697--759, 777--822, 1167--3755. Nothing inside a retained excerpt is omitted or altered: every retained body is delivered exactly as the article's lines are written, so no omission receipt of any kind accompanies this package. Citations: the retained excerpts carry no citation command at all, so no bibliography entry of the article is transcribed and no reference is redistributed. Reference adaptation: none was needed and none was made; every reference inside the delivered unit resolves inside it, so no unresolved reference or citation is produced. Printed slips kept literal: no printed word, symbol, hypothesis or number of the counted statement or of its proof was altered. Layout and class: the article's own class is \texttt{amsart} with packages including amsmath, amsthm, amscd, amsfonts, amssymb, xy, graphicx, enumerate, tikz, datetime; the wrapper is the lane's pinned \texttt{amsart} editorial wrapper at 10pt on A4, which restates the theorem-like environments the retained text needs and adds the article's own macro definitions that the retained text needs (none), with extra wrapper packages (bm, bbm, dsfont, xspace, cancel, mathtools). The delivered numbering is the unit's own. Assets: no retained span contains a figure environment, a graphics-inclusion command, a PDF-page inclusion, a table or a TikZ picture; no image file is copied into this package, the package declares no asset dependency and no image file is redistributed with it. Licence: version arXiv:2508.19835v1 of this article is distributed under the Creative Commons Attribution 4.0 International licence, as recorded in the arXiv licence metadata for this exact version and shown on the abstract page https://arxiv.org/abs/2508.19835v1. A full rights notice is delivered at licenses/arxiv-2508.19835-CC-BY-4.0.txt. Characters: every retained excerpt is pure ASCII, so the delivered engine, XeLaTeX with its default Latin Modern text font, reports no missing character.
\begin{proposition}
    \label{esfriou, vai dar tainha}
Let $(\mathcal{G}, X)$ be a relative ultragraph. Then there exists a 
$*$-homomorphism 
\[
\phi: C^*(\mathcal{G},X) \longrightarrow C^*(\mathcal{G}_X)
\]
such that, for each $B \in \mathcal{E}$,
\[
\phi(p_B) = P_B + P_{(B \cap Y)'}
\]
and, for each $e \in \mathcal{G}^1$,
\[
\phi(s_e) =
\begin{cases}
    S_e, & \text{if } r(e) \cap Y = \emptyset, \\[6pt]
    S_e + S_{e'}, & \text{if } r(e) \cap Y \neq \emptyset.
\end{cases}
\]
\end{proposition}

\begin{corollary}\label{decomposition2}
Let $(\mathcal{G},X)$ be a relative ultragraph and set $Y=\mathrm{Reg}(\mathcal{G})\setminus X$.
If $r(e)\cap Y$ is finite for every $e\in\mathcal{G}^1$, then for each $Z\in\mathcal{E}_X$ there
is a unique decomposition
\[
Z \;=\; A \;\cup\; (B\cap Y)',
\]
with $A,B\in\mathcal{E}$ and $(B\cap Y)'$ finite.
\end{corollary}

\begin{lemma}\label{contas1}
Let $(\mathcal{G},X)$ be a relative ultragraph and set $Y:=\mathrm{Reg}(\mathcal{G})\setminus X$.
Suppose $Z_1=A_1\cup (B_1\cap Y)'$ and $Z_2=A_2\cup (B_2\cap Y)'$ with $A_1,A_2,B_1,B_2\in\mathcal{E}$.
Define, for $w\in \mathcal{G}^0$,
\[
q_w \;:=\; p_w \;-\; \sum_{\{e\in \mathcal{G}^1:\ s(e)=w\}} s_e s_e^* .
\]
Then the following hold:
\begin{enumerate}
    \item \(p_{A_1} \left( \sum\limits_{e \mid s(e) \in A_2 \cap Y} s_e s_e^* \right) = p_{A_1 \cap Y} \left( \sum\limits_{e \mid s(e) \in A_2 \cap Y} s_e s_e^* \right)\).
    \item \(\left( \sum\limits_{e \mid s(e) \in A_1 \cap Y} s_e s_e^* \right) p_{A_2} = \left( \sum\limits_{e \mid s(e) \in A_1 \cap Y} s_e s_e^* \right) p_{A_2 \cap Y}\).
    \item \(p_{A_1} \left( \sum\limits_{w \mid w \in B_2 \cap Y} q_w \right) = p_{A_1 \cap Y} \left( \sum\limits_{w \mid w \in B_2 \cap Y} q_w \right)\).
    \item \(\left( \sum\limits_{w \mid w \in B_1 \cap Y} q_w \right) p_{A_2} = \left( \sum\limits_{w \mid w \in B_1 \cap Y} q_w \right) p_{A_2 \cap Y}\).
    \item \(\left( \sum\limits_{e \mid s(e) \in A_1 \cap Y} s_e s_e^* \right) \left( \sum\limits_{w \in B_2 \cap Y} q_w \right) = 0 = \left( \sum\limits_{w \in B_1 \cap Y} q_w \right) \left( \sum\limits_{e \mid s(e) \in A_2 \cap Y} s_e s_e^* \right)\).
    \item \(q_w q_v = \delta_{wv} q_w\), and consequently, \(\left( \sum\limits_{w \in B_1 \cap Y} q_w \right) \left( \sum\limits_{w \in B_2 \cap Y} q_w \right) = \sum\limits_{w \in B_1 \cap B_2 \cap Y} q_w\).
\end{enumerate}
\end{lemma}

\begin{proposition}\label{isomorphismrelative}
Let $(\mathcal{G},X)$ be a relative ultragraph and set
$Y:=\mathrm{Reg}(\mathcal{G})\setminus X$. If $r(e)\cap Y$ is finite for every
$e\in\mathcal{G}^1$, then the $*$-homomorphism
$\phi: C^*(\mathcal{G},X)\to C^*(\mathcal{G}_X)$ from
Proposition~\ref{esfriou, vai dar tainha} is injective.
\end{proposition}

\begin{proof}
Following the notation of Lemma~\ref{contas1}, define, for all
$w \in \mathrm{Reg}(\mathcal{G})=\mathrm{Reg}(\mathcal{G}_X)$,
\[
q_w \;:=\; p_w \;-\; \sum_{\substack{e \in \mathcal{G}^1 \\ s(e)=w}} s_e s_e^* .
\]
Let $Z \in \mathcal{E}_X^0$. By Corollary~\ref{decomposition2}, we can write $Z$ uniquely as
$Z = A \cup (B \cap Y)'$ where $A,B \in \mathcal{E}^0$ and $(B \cap Y)'$ is finite. We define $\psi$
on the generators of $C^*(\mathcal{G}_X)$ as follows:
\[
\psi(P_Z)
\;=\;
p_A - p_{A \cap Y}
\;+\;
\sum_{\substack{e \in \mathcal{G}^1 \\ s(e) \in A \cap Y}} s_e s_e^*
\;+\;
\sum_{\substack{w \in \mathcal{G}^0 \\ w \in B \cap Y}} q_w,
\]
and
\[
\psi(S_f)
=
\begin{cases}
\displaystyle
s_f \!\left(
p_{r(f)}
-
p_{r(f)\cap Y}
+
\sum_{\substack{g \in \mathcal{G}^1 \\ s(g) \in r(f)\cap Y}}
s_g s_g^*
\right),
& \text{if } f \in \mathcal{G}^1, \\[10pt]
\displaystyle
s_e \!\left(
p_{r(e)\cap Y}
-
\sum_{\substack{g \in \mathcal{G}^1 \\ s(g) \in r(e)\cap Y}}
s_g s_g^*
\right),
& \text{if } f=e' \in C.
\end{cases}
\]
Sums over the empty set are taken to be $0$. The sums written above are finite by the hypothesis.

We claim that $\psi$ is well defined and that $\psi \circ \phi = \mathrm{id}$. To see that $\psi$ is
well defined, we use the universal property of $C^*(\mathcal{G}_X)$; some computations are necessary:

\begin{enumerate}
\item We verify that $\psi(P_\emptyset)=0$, that
$\psi(P_{Z_1 \cap Z_2}) = \psi(P_{Z_1}) \psi(P_{Z_2})$, and that
$\psi(P_{Z_1\cup Z_2})= \psi(P_{Z_1})+\psi(P_{Z_2})-\psi(P_{Z_1\cap Z_2})$ for all
$Z_1, Z_2 \in \mathcal{E}_X^0$.

Of course $\psi(P_{\emptyset}) = p_{\emptyset} = 0$. Let us now prove that
$\psi(P_{Z_1 \cap Z_2}) = \psi(P_{Z_1}) \psi(P_{Z_2})$ for all $Z_1, Z_2 \in \mathcal{E}_X^0$.
Write
\[
Z_1 = A_1 \cup (B_1 \cap Y)'
\quad\text{and}\quad
Z_2 = A_2 \cup (B_2 \cap Y)'
\]
as in Corollary~\ref{decomposition2}. Note that $\psi(P_{Z_1}) \psi(P_{Z_2})$ is the product of
\[
\left(
p_{A_1} - p_{A_1 \cap Y}
+ \sum_{\substack{e \in \mathcal{G}^1 \\ s(e) \in A_1 \cap Y}} s_e s_e^*
+ \sum_{\substack{w \in \mathcal{G}^0 \\ w \in B_1 \cap Y}} q_w
\right)
\]
and
\[
\left(
p_{A_2} - p_{A_2 \cap Y}
+ \sum_{\substack{e \in \mathcal{G}^1 \\ s(e) \in A_2 \cap Y}} s_e s_e^*
+ \sum_{\substack{w \in \mathcal{G}^0 \\ w \in B_2 \cap Y}} q_w
\right).
\]
Using the equalities provided in Lemma~\ref{contas1}, most of the products are $0$ or cancel with
another term, so that
\[
\psi(P_{Z_1})\psi(P_{Z_2})
=
p_{A_1 \cap A_2}
-
p_{A_1 \cap A_2 \cap Y}
+
\sum_{\substack{e \in \mathcal{G}^1 \\ s(e) \in A_1 \cap A_2 \cap Y}} s_e s_e^*
+
\sum_{w \in B_1 \cap B_2 \cap Y} q_w.
\]
Thus we have the desired result; indeed, the right-hand side is exactly the definition of
$\psi(P_{Z_1 \cap Z_2})$, because
$Z_1 \cap Z_2 = (A_1 \cap A_2) \cup (B_1 \cap B_2 \cap Y)'$.

To finalize item~(1), we still need to show that
$\psi(P_{Z_1}) + \psi(P_{Z_2}) - \psi(P_{Z_1 \cap Z_2}) = \psi(P_{Z_1 \cup Z_2})$.
This is similar to what we did above, and we leave it to the interested reader.

\item In this item we show that $\psi(S_f)^*\psi(S_f) = \psi(P_{r(f)})$. If $f \in \mathcal{G}^1$,
then $\psi(S_f)^* \psi(S_f)$ equals
\[
\left(
s_f - s_f p_{r(f) \cap Y}
+ \sum_{\substack{g \in \mathcal{G}^1 \\ s(g) \in r(f) \cap Y}} s_f s_g s_g^*
\right)^{\!*}
\left(
s_f - s_f p_{r(f) \cap Y}
+ \sum_{\substack{g \in \mathcal{G}^1 \\ s(g) \in r(f) \cap Y}} s_f s_g s_g^*
\right).
\]
Using that, for all $g$ with $s(g) \in r(f) \cap Y$, we have
$p_{r(f) \cap Y}s_g s_g^* = s_g s_g^*$ and $p_{r(f)}s_g s_g^* = s_g s_g^*$, we conclude that
\[
\psi(S_f)^* \psi(S_f)
=
p_{r(f)} - p_{r(f) \cap Y}
+ \sum_{\substack{g \in \mathcal{G}^1 \\ s(g) \in r(f) \cap Y}} s_g s_g^*
=
\psi(P_{r(f)}),
\]
where the last equality holds because $r(f) \subseteq \mathcal{G}^0$.

The remaining case is when $f = e'$ for some $e \in \mathcal{G}^1$. In this situation,
$r(f) = r(e') = \emptyset \cup (r(e) \cap Y)'$, and thus
\[
\psi(P_{r(f)})
=
\psi(P_{r(e')})
=
\sum_{\substack{w \in \mathcal{G}^0 \\ w \in r(e) \cap Y}} q_w
=
p_{r(e) \cap Y}
-
\sum_{\substack{g \in \mathcal{G}^1 \\ s(g) \in r(e) \cap Y}} s_g s_g^*,
\]
where the last equality holds because the right-hand side is the left-hand side after reordering
the sum. Moreover, a computation shows that
\begin{align*}
\psi(S_f)^*\psi(S_f)
&=
\left(
s_e  p_{r(e) \cap Y}
-
\sum_{\substack{g \in \mathcal{G}^1 \\ s(g) \in r(e) \cap Y}} s_e s_g s_g^*
\right)^{\!*}
\left(
s_e  p_{r(e) \cap Y}
-
\sum_{\substack{g \in \mathcal{G}^1 \\ s(g) \in r(e) \cap Y}} s_e s_g s_g^*
\right) \\
&=
\left(
p_{r(e) \cap Y} s_e^*
-
\sum_{\substack{g \in \mathcal{G}^1 \\ s(g) \in r(e) \cap Y}} s_g s_g^* s_e^*
\right)
\left(
s_e  p_{r(e) \cap Y}
-
\sum_{\substack{g \in \mathcal{G}^1 \\ s(g) \in r(e) \cap Y}} s_e s_g s_g^*
\right) \\
&=
p_{r(e) \cap Y}
-
\sum_{\substack{g \in \mathcal{G}^1 \\ s(g) \in r(e) \cap Y}} s_g s_g^*
-
\sum_{\substack{g \in \mathcal{G}^1 \\ s(g) \in r(e) \cap Y}} s_g s_g^*
+
\sum_{\substack{g \in \mathcal{G}^1 \\ s(g) \in r(e) \cap Y}} s_g s_g^* \\
&=
p_{r(e) \cap Y}
-
\sum_{\substack{g \in \mathcal{G}^1 \\ s(g) \in r(e) \cap Y}} s_g s_g^*
=
\psi(P_{r(f)}).
\end{align*}

\item We need to prove that, for all $f \in \mathcal{G}_X^1$,
\[
\psi(P_{s(f)}) \,\psi(S_f)\, \psi(S_f)^* \;=\; \psi(S_f)\, \psi(S_f)^*.
\]
To start, let us compute $\psi(S_f) \psi(S_f)^*$; this computation will be useful also in the next
item. There are two possible situations:

\begin{itemize}
\item \textbf{Case 1}. $f \in \mathcal{G}^1$. Using the definition of $\psi$ and performing some computations, we obtain
\begin{align*}
\psi(S_f)\psi(S_f)^*
&= s_f s_f^* - s_f\, p_{r(f)\cap Y}\, s_f^*
  + \sum_{\substack{g\in\mathcal{G}^1\\ s(g)\in r(f)\cap Y}}
    s_f s_g s_g^* s_f^* \\[2pt]
&= s_f\!\left(
    s_f^* - p_{r(f)\cap Y}\, s_f^*
    + \sum_{\substack{g\in\mathcal{G}^1\\ s(g)\in r(f)\cap Y}}
      s_g s_g^* s_f^*
  \right).
\end{align*}

\item \textbf{Case 2}. $f = e' \in C$. Using the definition of $\psi$ and performing some computations, we obtain
\[
\psi(S_{e'})\psi(S_{e'})^*
=
s_e p_{r(e) \cap Y} s_e^*
-
\sum_{\substack{g \in \mathcal{G}^1 \\ s(g) \in r(e) \cap Y}} s_e s_g s_g^* s_e^*
=
s_e \!\left(
p_{r(e) \cap Y} s_e^*
-
\sum_{\substack{g \in \mathcal{G}^1 \\ s(g) \in r(e) \cap Y}} s_g s_g^* s_e^*
\right)\!.
\]
\end{itemize}

In both cases, we now compute the product $\psi(P_{s(f)})\,\psi(S_f)\,\psi(S_f)^*$. Note that, for all $f \in \mathcal{G}_X^1$,
$s(f) \in \mathcal{G}^0$; moreover,
\[
\psi(P_{s(f)}) \;=\;
\begin{cases}
p_{s(f)}, & \text{if } s(f) \notin Y, \\[2pt]
\displaystyle \sum_{\substack{e \in \mathcal{G}^1 \\ s(e) = s(f)}} s_e s_e^*, & \text{if } s(f) \in Y.
\end{cases}
\]
Thus, in each case above we need to consider two subcases. The reader can check all these
possibilities by direct computations; they are very similar and use the facts that
$p_{s(f)} s_f = s_f$ and
$\sum_{\substack{e \in \mathcal{G}^1 \\ s(e) = s(f)}} s_e s_e^*\, s_f = s_f$.

The three cases above can be used to check that the $\psi(S_f)$ are partial isometries and that
$\psi(S_f) \psi(S_f)^* \, \psi(S_g) \psi(S_g)^* = 0$ if $f \neq g$. We leave this to the reader.

\item In this item we show that, for all $v \in \mathrm{Reg}(\mathcal{G}_X)=\mathrm{Reg}(\mathcal{G})$,
\[
\psi(P_v) \;=\; \sum_{\substack{f \in \mathcal{G}_X^1 \\ s(f) = v}} \psi(S_f)\, \psi(S_f)^*.
\]
For notational convenience, set $O = \{\,f \in \mathcal{G}^1 \mid r(f) \cap Y = \emptyset\,\}$ and
$Q = \{\,f \in \mathcal{G}^1 \mid r(f) \cap Y \neq \emptyset\,\}$ only for this item. Note that
$\mathcal{G}^1 = O \cup Q$, thus
\[
\mathcal{G}_X^1 = O \cup Q \cup C.
\]
Using the computations made in item~(3), we note that
\[
\sum_{\substack{f \in O \\ s(f) = v}} \psi(S_f) \psi(S_f)^*
=
\sum_{\substack{f \in O \\ s(f) = v}} s_f s_f^*,
\]
\[
\sum_{\substack{f \in Q \\ s(f) = v}} \psi(S_f) \psi(S_f)^*
=
\sum_{\substack{f \in Q \\ s(f) = v}}
\Big(
s_f s_f^*
-
s_f p_{r(f) \cap Y} s_f^*
+
\sum_{\substack{g \in \mathcal{G}^1 \\ s(g) \in r(f) \cap Y}}
s_f s_g s_g^* s_f^*
\Big),
\]
and
\[
\sum_{\substack{f' \in C \\ s(f') = v}} \psi(S_{f'}) \psi(S_{f'})^*
=
\sum_{\substack{f' \in C \\ s(f') = v}}
\Big(
s_f p_{r(f) \cap Y} s_f^*
-
\sum_{\substack{g \in \mathcal{G}^1 \\ s(g) \in r(f) \cap Y}}
s_f s_g s_g^* s_f^*
\Big).
\]
Using that $s(f') = s(f)$ for all $f \in \mathcal{G}^1$ with $r(f) \cap Y \neq \emptyset$, we can sum
the three equalities above to obtain
\begin{align*}
\sum_{\substack{f \in \mathcal{G}_X^1 \\ s(f) = v}} \psi(S_f) \psi(S_f)^*
&=
\sum_{\substack{f \in O \\ s(f) = v}} s_f s_f^*
\;+\;
\sum_{\substack{f \in Q \\ s(f) = v}} s_f s_f^* \\
&=
\sum_{\substack{f \in \mathcal{G}^1 \\ s(f) = v}} s_f s_f^*.
\end{align*}
Since $v \in \mathcal{G}^0$ and $\{v\} = \{v\} \cup \emptyset$, we have two options: $v \notin Y$ or
$v \in Y$. In the first case, $\psi(P_v) = p_v$ and, since $v \notin Y$ (i.e., $v \in X$),
$p_v = \sum_{\substack{f \in \mathcal{G}^1 \\ s(f) = v}} s_f s_f^*$. In the second case, when
$v \in Y$, we have $\psi(P_v) = \sum_{\substack{f \in \mathcal{G}^1 \\ s(f) = v}} s_f s_f^*$.
In both cases we obtain the desired result. With these four items, we conclude that $\psi$ extends
to a well-defined homomorphism.

To conclude the theorem, it remains to check that $\psi \circ \phi = \mathrm{id}_{C^*(\mathcal{G},X)}$.
First, for all $v \in \mathcal{G}^0$ we have $\psi \circ \phi(p_v) = p_v$; this is immediate if
$v \notin Y$, because in this case $\psi \circ \phi(p_v) = \psi(P_v) = p_v$. If $v \in Y$, then
\[
\psi \circ \phi(p_v)
=
\psi(P_v + P_{v'})
=
\sum_{\substack{e \in \mathcal{G}^1 \\ s(e) = v}} s_e s_e^*
\;+\;
\left( p_v - \sum_{\substack{e \in \mathcal{G}^1 \\ s(e) = v}} s_e s_e^* \right)
=
p_v.
\]
Next, for all $e \in \mathcal{G}^1$,
\begin{align*}
\psi \circ \phi(p_{r(e)})
&= \psi\big( P_{r(e)} + P_{(r(e) \cap Y)'} \big) \\
&= p_{r(e)} - p_{r(e) \cap Y}
+ \sum_{\substack{g \in \mathcal{G}^1 \\ s(g) \in r(e) \cap Y}} s_g s_g^*
+ \sum_{w \in r(e) \cap Y} q_w \\
&= p_{r(e)} - p_{r(e) \cap Y}
+ \sum_{\substack{g \in \mathcal{G}^1 \\ s(g) \in r(e) \cap Y}} s_g s_g^*
+ p_{r(e) \cap Y}
- \sum_{\substack{g \in \mathcal{G}^1 \\ s(g) \in r(e) \cap Y}} s_g s_g^* \\
&= p_{r(e)}.
\end{align*}
Finally, let $e \in \mathcal{G}^1$. If $r(e) \cap Y = \emptyset$, then
\[
\psi \circ \phi(s_e) = \psi(S_e) = s_e.
\]
If $r(e) \cap Y \neq \emptyset$, then
\begin{align*}
\psi \circ \phi(s_e)
&= \psi(S_e + S_{e'}) \\
&= s_e - s_e p_{r(e) \cap Y}
+ \sum_{\substack{g \in \mathcal{G}^1 \\ s(g) \in r(e) \cap Y}} s_e s_g s_g^*
\;+\;
s_e p_{r(e) \cap Y}
- \sum_{\substack{g \in \mathcal{G}^1 \\ s(g) \in r(e) \cap Y}} s_e s_g s_g^* \\
&= s_e.
\end{align*}
Since $\psi \circ \phi = \mathrm{id}_{C^*(\mathcal{G},X)}$ on the set of generators of
$C^*(\mathcal{G},X)$, the result follows.
\end{enumerate}
\end{proof}

\end{document}
